\documentclass[12pt, a4paper]{article}

\usepackage{amsmath}
\usepackage{amssymb}
\usepackage{booktabs}
\usepackage{graphicx}
\usepackage{caption}
\usepackage{geometry}
\usepackage{hyperref}
\usepackage{siunitx}
\usepackage{array}
\usepackage{float}
\usepackage{setspace}
\usepackage{circuitikz}
\usetikzlibrary{positioning}

\onehalfspacing

\geometry{margin=2.5cm}

\hypersetup{
    colorlinks=true,
    linkcolor=blue,
    urlcolor=blue,
    citecolor=blue
}

\sisetup{per-mode=symbol}

\title{A Discrete-Component Second-Order CT Sigma-Delta ADC Design}
\author{}
\date{}

\begin{document}

\maketitle
\tableofcontents
\newpage

% ─────────────────────────────────────────────
\section{Overview}
\label{section:overview}

In the sigma-delta ($\Delta\Sigma$) ADC topology, a quantizer produces a stream of coarse 1-bit decisions, each of which is fed back through a 1-bit DAC into the loop's cascaded integrators, shaping the input to subsequent decisions. Because resolution is built up from many such decisions rather than a single high-precision comparison, a $\Delta\Sigma$ ADC trades conversion latency for resolution relative to a SAR ADC, which resolves a full-precision word in a single conversion cycle.
This report documents a second-order continuous-time sigma-delta ADC, designed and tested in
LTspice before discrete breadboard implementation using physical components. The target specifications
for the design are listed in Table~\ref{tab:specs}.

\begin{table}[H]
    \centering
    \caption{Design Specifications}
    \label{tab:specs}
    \begin{tabular}{@{}ll@{}}
        \toprule
        Parameter & Value \\
        \midrule
        Signal Bandwidth & \qty{200}{\hertz} \\
        Oversampling Ratio (OSR) & 128 \\
        Sample Rate & \qty{51.2}{\kilo\hertz} \\
        \bottomrule
    \end{tabular}
\end{table}

% ─────────────────────────────────────────────
\section{Architecture}
\label{section:architecture}



% ─────────────────────────────────────────────
\section{Bring-Up Methodology}
\label{section:bring_up}

% ------------------------------------------------------------
\subsection{Root Cause: DAC Sign Error}
\label{section:dac_sign_error}

% ------------------------------------------------------------
\subsection{Debugging Highlight: Instrumentation Artefacts}
\label{section:instrumentation_artefacts}

% ------------------------------------------------------------
\section{Results}
\label{section:results}

% ------------------------------------------------------------
\section{Limitations}
\label{section:limitations}

% ------------------------------------------------------------
\section{Next Steps}
\label{section:next_steps}

% ------------------------------------------------------------
\addcontentsline{toc}{section}{References}
\begin{thebibliography}{9}
\end{thebibliography}

% ------------------------------------------------------------
\appendix

\section{Source Code and Reproducibility}
\label{appendix:source}

\end{document}